%% Font size %%
\documentclass[11pt]{article}

%% Load the custom package
\usepackage{Mathdoc}

%% Numéro de séquence %% Titre de la séquence %%
\renewcommand{\centerhead}{Chapitre 01 : Nombres relatifs - Soustractions}

%% Spacing commands %%
\renewcommand{\baselinestretch}{1} \setlength{\parindent}{0pt}

\begin{document}

\begin{exercice}[1]
Effectuer les calculs suivants comme dans l'exemple.

\begin{multicols}{2}

$A = (-7) - (-9)$\\
$= (-7) + (+9)$\\
$= +(9 - 7)$\\
$= +2$

\bigskip

$B = (+11) - (-14)$\\
$= \dotfill$\\
$= \dotfill$\\
$= \dotfill$

\bigskip

$C = (-16) - (+5)$\\
$= \dotfill$\\
$= \dotfill$\\
$= \dotfill$

\columnbreak

$D = (+8) - (+19)$\\
$= \dotfill$\\
$= \dotfill$\\
$= \dotfill$

\bigskip

$E = (-3) - (-3)$\\
$= \dotfill$\\
$= \dotfill$\\
$= \dotfill$

\bigskip

$F = (-12) - (+17)$\\
$= \dotfill$\\
$= \dotfill$\\
$= \dotfill$

\end{multicols}
\end{exercice}

\begin{exercice}[2]
Effectuer les calculs suivants comme dans l'exemple.

\begin{multicols}{2}

$A = (-3{,}5) - (-2{,}1)$\\
$= (-3{,}5) + (+2{,}1)$\\
$= -(3{,}5 - 2{,}1)$\\
$= -1{,}4$

\bigskip

$B = (+7{,}3) - (+4{,}8)$\\
$= \dotfill$\\
$= \dotfill$\\
$= \dotfill$

\bigskip

$C = (-2{,}6) - (+5{,}1)$\\
$= \dotfill$\\
$= \dotfill$\\
$= \dotfill$

\columnbreak

$D = (+0{,}9) - (-6{,}4)$\\
$= \dotfill$\\
$= \dotfill$\\
$= \dotfill$

\bigskip

$E = (-8{,}2) - (-1{,}7)$\\
$= \dotfill$\\
$= \dotfill$\\
$= \dotfill$

\bigskip

$F = (-5{,}4) - (+2{,}7)$\\
$= \dotfill$\\
$= \dotfill$\\
$= \dotfill$

\end{multicols}
\end{exercice}

\newpage

\begin{exercice}[2][Compléter]
Trouver le nombre manquant dans chaque égalité, comme dans l'exemple.

\begin{exemple}
$(+15) - \ldots\ldots\ldots = (+8)$\\
On cherche : $(+15) - (+8) = +7$, donc le nombre manquant est $+7$.
\end{exemple}

\begin{multicols}{2}
\begin{enumerate}[itemsep=1.5em]
\item $(-9) - \ldots\ldots\ldots = (-14)$
\item $\ldots\ldots\ldots - (+7) = (-3)$
\item $(-2) - \ldots\ldots\ldots = (+6)$
\item $\ldots\ldots\ldots - (-5) = (-1)$
\item $(+9) - \ldots\ldots\ldots = (+13)$
\end{enumerate}
\end{multicols}
\end{exercice}

\begin{exercice}[3]
Effectuer les calculs suivants comme dans l'exemple.

\begin{multicols}{2}

$A = (-62) - (-85)$\\
$= (-62) + (+85)$\\
$= +(85 - 62)$\\
$= +23$

\bigskip

$B = (+77) - (+94)$\\
$= \dotfill$\\
$= \dotfill$\\
$= \dotfill$

\bigskip

$C = (-58) - (+63)$\\
$= \dotfill$\\
$= \dotfill$\\
$= \dotfill$

\columnbreak

$D = (+91) - (-52)$\\
$= \dotfill$\\
$= \dotfill$\\
$= \dotfill$

\bigskip

$E = (-99) - (-70)$\\
$= \dotfill$\\
$= \dotfill$\\
$= \dotfill$

\bigskip

$F = (-73) - (+56)$\\
$= \dotfill$\\
$= \dotfill$\\
$= \dotfill$

\end{multicols}
\end{exercice}

\end{document}
